% from ega4-2-fr.tex
\oldpage[IV-2]{57}

\item[(iii)] Les indices $\lambda$ tels que $x\in Z_\lambda$ sont les mêmes
que ceux pour lesquels il existe un indice $\mu$ tel que
$x'\in Z_{\lambda\mu}$. En particulier, si $x'$ n'appartient qu'à un seul
cycle premier associé à $\mathcal F\otimes_k K$, $x$ n'appartient qu'à un
seul cycle premier associé à $\mathcal F$.
\item[(iv)] Si $X$ est localement de type fini sur $k$, on a
$\dim(Z_{\lambda\mu})=\dim(Z_\lambda)$.
\end{enumerate}
\end{proposition}

\begin{proof}
On notera que l'hypothèse entraîne que $X$ lui-même est localement noethérien
(2.2.13 et 2.2.14)~; l'assertion (i) résulte de (4.2.5) et de la démonstration
de (4.2.3), pour $\mathcal G=\mathcal O_Y$, avec
$Y=\operatorname{Spec}(K)$~; (ii) et (iii) résultent de (i) et de (2.3.5)~;
enfin, (iv) est un cas particulier de (4.2.4, (iii)).
\end{proof}

\begin{corollary}[4.2.8]
\label{IV.2.4.2.8-fr}
Si $X$ est localement de type fini sur $k$, l'ensemble des dimensions des
cycles premiers associés est le même pour $\mathcal F$ et
$\mathcal F\otimes_k K$~; l'ensemble des dimensions des composantes
irréductibles est le même pour $X$ et $X\otimes_k K$.
\end{corollary}

\begin{proposition}[4.2.9]
\label{IV.2.4.2.9-fr}
Supposons vérifiées les hypothèses de (4.2.5), et supposons en outre que
$\mathcal F$ et $\mathcal G$ soient cohérents. Soient
$(\mathcal F_\lambda)_{\lambda\in L}$ et
$(\mathcal G_\mu)_{\mu\in M}$ des décompositions irredondantes de
$\mathcal F$ et $\mathcal G$ respectivement~; pour tout couple
$(\lambda,\mu)\in L\times M$, soit
$(\mathcal H_{\lambda\mu\nu})_{\nu\in S(\lambda,\mu)}$ une décomposition
irredondante réduite de $\mathcal F_\lambda\otimes_k\mathcal G_\mu$, où
$S(\lambda,\mu)=\operatorname{Ass}(\mathcal F_\lambda\otimes_k\mathcal G_\mu)$
(3.2.5). Alors $(\mathcal H_{\lambda\mu\nu})$ (pour tous les triplets
$(\lambda,\mu,\nu)$) est une décomposition irredondante de
$\mathcal F\otimes_k\mathcal G$~; elle est réduite si
$(\mathcal F_\lambda)$ et $(\mathcal G_\mu)$ le sont.
\end{proposition}

\begin{proof}
Notons que $\mathcal F\otimes_k\mathcal G$ et chacun des
$\mathcal F_\lambda\otimes_k\mathcal G_\mu$ sont cohérents, et chacun des
$\mathcal F_\lambda\otimes_k\mathcal G_\mu$ s'identifie à un quotient de
$\mathcal F\otimes_k\mathcal G$~; chacun des $\mathcal H_{\lambda\mu\nu}$
s'identifie par définition à un quotient cohérent de
$\mathcal F_\lambda\otimes_k\mathcal G_\mu$, donc aussi à un quotient
cohérent de $\mathcal F\otimes_k\mathcal G$. La famille des supports des
$\mathcal H_{\lambda\mu\nu}$ est localement finie, car
$\operatorname{Supp}(\mathcal F_\lambda\otimes_k\mathcal G_\mu)$ est contenu
dans l'espace sous-jacent à
$\operatorname{Supp}(\mathcal F_\lambda)\times_k
\operatorname{Supp}(\mathcal G_\mu)$ (où
$\operatorname{Supp}(\mathcal F_\lambda)$ et
$\operatorname{Supp}(\mathcal G_\mu)$ désignent les sous-préschémas fermés
réduits correspondants), et pour $\lambda$ et $\mu$ donnés, la famille des
supports des $\mathcal H_{\lambda\mu\nu}$ est localement finie par hypothèse~;
notre assertion résulte donc de ce que les familles
$(\operatorname{Supp}(\mathcal F_\lambda))$ et
$(\operatorname{Supp}(\mathcal G_\mu))$ sont localement finies. Pour démontrer
que $(\mathcal H_{\lambda\mu\nu})$ est une décomposition irredondante de
$\mathcal F\otimes_k\mathcal G$, il suffit donc de prouver que
l'homomorphisme canonique
\[
\mathcal F\otimes_k\mathcal G
\longrightarrow\bigoplus_{\lambda,\mu,\nu}\mathcal H_{\lambda\mu\nu}
\]
est injectif~; or, il est composé des homomorphismes
\[
\mathcal F\otimes_k\mathcal G
\longrightarrow
\left(\bigoplus_\lambda\mathcal F_\lambda\right)\otimes_k
\left(\bigoplus_\mu\mathcal G_\mu\right)
\longrightarrow\bigoplus_{\lambda,\mu,\nu}\mathcal H_{\lambda\mu\nu}~;
\]
le second de ces homomorphismes est injectif par définition, et il en est de
même du premier, qui n'est autre que le produit tensoriel des homomorphismes
canoniques injectifs
$\mathcal F\to\bigoplus_\lambda\mathcal F_\lambda$,
$\mathcal G\to\bigoplus_\mu\mathcal G_\mu$ (on se souviendra que
$\mathcal F$ et $\mathcal G$ sont \emph{plats sur $k$}). Enfin, si
$(\mathcal F_\lambda)$ et $(\mathcal G_\mu)$ sont réduites, on peut supposer
que $L=\operatorname{Ass}(\mathcal F)$ et
$M=\operatorname{Ass}(\mathcal G)$~; le fait que les
$\mathcal H_{\lambda\mu\nu}$ forment alors une décomposition réduite résulte
de ce que les $\operatorname{Ass}(\mathcal F_\lambda\otimes_k
\mathcal G_\mu)$ forment une partition de
$\operatorname{Ass}(\mathcal F\otimes\mathcal G)$ en vertu de (4.2.2)
(cf. (3.2.5)).
\end{proof}

\begin{corollary}[4.2.10]
\label{IV.2.4.2.10-fr}
Sous les hypothèses de (4.2.7), supposons en outre $\mathcal F$ cohérent, et
soit $(\mathcal F_\lambda)_{\lambda\in L}$ une décomposition irredondante de
$\mathcal F$~; pour tout $\lambda\in L$, soit
$(\mathcal F_{\lambda\mu})_{\mu\in
\operatorname{Ass}(\mathcal F_\lambda\otimes_k K)}$ une décomposition
irredondante réduite de $\mathcal F_\lambda\otimes_k K$. Alors
$(\mathcal F_{\lambda\mu})$ est une décomposition irredondante de
$\mathcal F\otimes_k K$~; elle est réduite si $(\mathcal F_\lambda)$ l'est.
\end{corollary}

\begin{proof}
On applique (4.2.9) avec $Y=\operatorname{Spec}(K)$,
$\mathcal G=\mathcal O_Y=K$, $M$ réduit à un seul élément.
\end{proof}

